\section{Worked families: what the parameters reveal}
\label{sec:worked-coverage}

The following examples distinguish the guarantees of compilation methods.
Their projected functions are explicitly easy. We calculate their actual
complexities where possible, so that a large upper-bound expression is
never mistaken for a lower bound on the function.

\subsection{An exponential improvement from selecting and sharing together}

Let $N\ge2$, $k\ge3$, and $\ell\ge2$. The ordinary inputs are
$z\in\B^N$, $u\in\B^k$; the witnesses are $a\in\B^k$,
$b\in\B^\ell$. Indices on $a$ are cyclic. Define
\begin{align}
 T(u,a)&=\AND_{i=1}^k[u_i=a_i\oplus a_{i+1}],\label{eq:cyclic-check}\\
 f(z,u;a,b)&=\operatorname{PAR}_N(z)\land T(u,a)
                         \land\operatorname{PAR}_\ell(b).
                         \label{eq:cyclic-verifier}
\end{align}
Use XOR chains for the two parities, one XOR and one XNOR for each
equation, a prefix AND chain for the equations, then AND the $z$ parity
with $T$ and finally with the $b$ parity. This supplied circuit has
\[
 n=N+k,\qquad m=k+\ell,\qquad s=N+3k+\ell-1.
\]
Every gate is in the output cone. Every input is essential: starting
from an accepting assignment, changing a $z$ or $b$ bit flips a required
parity; changing a $u$ bit violates an equation; changing an $a$ bit
violates two equations. Therefore
\[
 N+2k+\ell-1\le\C(f)\le N+3k+\ell-1.
\]
The supplied size is within $k$ gates of the elementary lower bound;
there is no padding or dead subcircuit.

XORing the equations shows that $T$ is satisfiable only if
$\operatorname{PAR}_k(u)=0$. Conversely, a choice of $a_1$ determines
all other $a$ bits along the first $k-1$ equations, and the last holds
precisely when that parity is zero. Hence
\begin{equation}\label{eq:cyclic-projection}
 g(z,u)=\operatorname{PAR}_N(z)\land\neg\operatorname{PAR}_k(u),
 \qquad \C(g)=N+k-1.
\end{equation}
The upper bound uses two parity chains and one binary gate; the lower
bound follows because all $N+k$ ordinary inputs are essential.

Each positive input has two complementary choices of $a$ and
$2^{\ell-1}$ odd-parity choices of $b$, giving exactly $2^\ell$
witnesses. The fibers for distinct even-parity $u$ are disjoint. Fix
one odd-parity $b_0$ and select
\[
 H=\{(a,b_0):a_1=0\},\qquad |H|=2^{k-1}=:h.
\]
Exactly one member of $H$ works for each even-parity $u$, so
$\tau(f)=h$. With
$A(H)=\sum_v N_v(H)+N_{\mathrm{out}}(H)-1$, the precise support
counts for the supplied circuit are as follows.

\begin{center}\small
\begin{tabular}{@{}lrr@{}}
\toprule
Circuit part & Selected $H$ copies & All support assignments\\
\midrule
$z$ parity & $N-1$ & $N-1$\\
$k$ XORs and $k$ XNORs & $8k-8$ & $8k$\\
Prefix conjunction & $3\cdot2^{k-1}-4$ & $3\cdot2^k-8$\\
$b$ parity & $\ell-1$ & $2^{\ell+1}-4$\\
AND of $z$ parity and $T$ & $2^{k-1}$ & $2^k$\\
Final AND & $2^{k-1}$ & $2^{k+\ell}$\\
Output OR & $2^{k-1}-1$ & $2^{k+\ell}-1$\\
\bottomrule
\end{tabular}
\end{center}

For the second row, two cycle edges touch $a_1$. Their gates each have
two selected assignments, while the remaining edges each have four.
For the third row, the $j$th prefix, $2\le j<k$, uses
$a_1,\ldots,a_{j+1}$; the last uses all $a$ bits. Thus its selected
sum is $\sum_{j=2}^{k-1}2^j+2^{k-1}$, and the full sum doubles it.
The parity-chain counts sum $2^2,\ldots,2^\ell$. Adding the rows gives
\begin{align*}
 A(H)&=N+\ell+8k+3\cdot2^k-15,\\
 A(\text{full})&=N+8k+4\cdot2^k+2^{\ell+1}+2^{k+\ell+1}-14.
\end{align*}
Exactly the $k$ witness vertices $a_i$ have fan-out two; all gate
outputs have fan-out one. Hence the original shared-vertex count is
$r=k$, and $s+1-n-m=k$.

Now set $N=2^k$ and $\ell=k$. The comparison becomes
\begin{center}
\begin{tabular}{@{}lr@{}}
\toprule
Quantity & Order as $k\to\infty$\\
\midrule
Actual $\C(g)$ and verifier size & $\Theta(2^k)$\\
Selected and shared cover $A(H)$ & $\Theta(2^k)$\\
Independent optimal-cover estimate $h(s+1)-1$ & $\Theta(4^k)$\\
Full-support estimate $A(\text{full})$ & $\Theta(4^k)$\\
Shared-vertex estimate $(s+1)2^r$ & $\Theta(4^k)$\\
Uniform-density and exhaustive estimates & $\Theta(8^k)$\\
\bottomrule
\end{tabular}
\end{center}
The universal synthesis ceiling is still larger, and the original
three-way exponent has $\min\{n,m,s+1-n-m\}=k$. Thus simultaneous
selection and sharing improve the minimum of all those original
guarantee expressions by a factor $\Theta(2^k)$. This does not say
that running an old compiler and then semantically simplifying it
requires $\Theta(4^k)$ gates. The new consistency-graph theorem can
also recognize this family's easy projection.

\subsection{Even every row mass leaves an exponential ambiguity}
\label{sec:row-mass-gap}

For the same diagonal family, let $\mu$ be uniform on $H$. Every
positive row has mass $1/h$, and
\[
 E_\mu(t)=2^{n-2}(1-1/h)^t,\qquad B_\mu(t)\le A(H)+1.
\]
Taking $t=\lceil h((n-2)\ln2+\ln n+1)\rceil$ makes
$nE_\mu(t)\le e^{-1}$. Although $t=\Theta(4^k)$, the bound in
\eqref{eq:occupancy-linear} is $O(2^k)$: repeated local restrictions
do not add gates.

Instead let $\nu$ be uniform on all witnesses whose $b$ parity is odd.
\emph{Every positive row still has exactly the same mass $1/h$.}
Both distributions maximize the minimum row mass, since the $h$
distinct positive fibers are disjoint. Nevertheless
\begin{equation}\label{eq:row-mass-gap}
 \inf_{t\ge0}\{B_\mu(t)+nE_\mu(t)\}=\Theta(2^k),
 \qquad
 \inf_{t\ge0}\{B_\nu(t)+nE_\nu(t)\}=\Theta(4^k),
\end{equation}
where the infima are over integer sample counts.

Here is a direct proof of the second assertion. There are
$M=2^{2k-1}$ equally likely witnesses under $\nu$, all distinguished
by the output support. Thus
$B_\nu(t)\ge M(1-(1-1/M)^t)$. If $t\ge M/8$, this is at least
$M(1-e^{-1/8})$. If $t<M/8$, use
$\ln(1-p)\ge-p/(1-p)$ with $p=1/h\le1/4$ to obtain
\[
 E_\nu(t)\ge2^{n-2}e^{-N/6}.
\]
This eventually exceeds every constant multiple of $4^k$, since
$n=N+k$ and $N=2^k$. These cases prove the lower bound on the
\emph{guarantee expression}. Its upper bound follows from the full
support count and $E_\nu(t)\to0$. The first assertion follows from
the preceding choice of $t$ and the actual lower bound on $\C(g)$.
The same exponential distinction survives the sparse correction
term, because $E/\log(E+2)$ in the second case is still exponentially
large in $N$.

Thus even the whole vector $(\mu(W_x))_{x:g(x)=1}$ does not determine
the cost of this compiler. How the same witnesses collide on internal
supports is additional information.

\subsection{Better coverage can be a worse circuit choice}

For $L\ge8$, put $n=2L$, $m=2$, and build
\[
 f_L(x;y)=
 (y_1\oplus x_1\oplus\cdots\oplus x_L)
 \lor
 (y_2\oplus x_{L+1}\oplus\cdots\oplus x_{2L}),
\]
with each XOR chain starting at its witness bit. There are $L$ gates
with each singleton witness support and one with both witnesses.
All inputs are essential, proving $\C(f_L)=2L+1$ by the input bound.
Its projection is nonetheless $g=1$.

Each row excludes just one of the four witnesses, and each excluded
label occurs equally often. Therefore
$\min_x\lambda(W_x)=1-\max_y\lambda(y)$, whose unique maximizing
distribution is the uniform distribution $u$, with value $3/4$.
For $F_\lambda(t)=B_\lambda(t)+nE_\lambda(t)$, direct occupancy gives
\begin{align*}
 B_u(t)&=4L(1-2^{-t})+8(1-(3/4)^t),&
 E_u(t)&=2^{2L-2t}.
\end{align*}
For $t<L$ the correction is at least $8L$; for $t\ge L$ the
occupancy is at least $4L(1-2^{-L})$. Hence
$F_u(t)\ge4L(1-2^{-L})$ for every integer $t\ge0$.

Let $v$ instead give equal mass to $00$ and $01$. Its minimum row
mass is only $1/2$. For $t\ge1$,
\[
 B_v(t)=L+(2L+4)(1-2^{-t}),\qquad E_v(t)=2^{2L-1-t}.
\]
At $t=2L+\lceil\log(2L)\rceil$ this gives
\[
 F_v(t)\le3L+4.5<4L(1-2^{-L})\le\inf_j F_u(j).
\]
The unique coverage optimizer is therefore worse for the circuit-cost
objective. The asymptotic comparison is $3L+O(1)$ versus $4L+O(1)$.
It also survives any fixed positive constant multiplying the sparse
correction: for $t<L/2$ that correction grows faster than $L$, while
for $t\ge L/2$ the uniform occupancy is already $4L-o(1)$.

This objective belongs to the supplied circuit representation. Moving
the witnesses to the ends of the XOR chains changes its supports
without changing $f_L$ or the optimal verifier size. Its constant
projection also emphasizes that an obstruction to a particular
compilation strategy is not a circuit lower bound.

\subsection{Two additional cautions}

Fiber sizes do not determine covers. For ordinary $u\in\B^k$ and
witnesses $(c,v)\in\B\times\B^k$, consider rows
\[
 W^A_u=\{(0,0^k),(1,u)\},\qquad
 W^Q_u=\{(0,u),(1,\neg u)\}.
\]
Both profiles consist entirely of size-two fibers; both have linear-size
verifiers with all inputs essential. The first has a common witness
and cover number one; the second partitions the witness universe
and has cover number $2^k$.

Nor may each gate choose its own best cover. In the relation
$f(x_1,x_2;y_1,y_2)=[x\ne y]$, any two distinct witnesses cover every
input. The set $\{00,01\}$ uses only one $y_1$ value, while
$\{00,10\}$ uses only one $y_2$ value. No cover achieves both
simultaneously, because it would fix both bits. The minimum of the
sum of these two occupancies is three, whereas the sum of their
separate minima is two. The occupancy theorem uses one common
distribution and one common sample throughout the circuit.
